\chapter{加法入门}\label{ch:g1:addition}

3 个红色弹珠和 2 个蓝色弹珠合在一起，一共有几个？
把它们合在一起并数出总数，就是\emph{加法}——第一种运算，
它有两个著名的符号 $+$ 和 $=$。

\section{合起来}

\begin{definition}[加法]\label{def:g1:addition:def}
\emph{加法}是把几部分合在一起并数出总数。我们
记作
\[
3 + 2 = 5
\]
读作：“三加二等于五”。加法的结果称为
\emph{和}\index{和}。
\end{definition}

\begin{omfigure}
\begin{tikzpicture}[scale=0.6]
  \foreach \i in {0,1,2} \fill[omThm] (\i*1.0,0) circle (0.3);
  \node[font=\large] at (3.0,0) {$+$};
  \foreach \i in {0,1} \fill[omDef] (3.8+\i*1.0,0) circle (0.3);
  \node[font=\large] at (5.8,0) {$=$};
  \foreach \i in {0,1,2} \fill[omThm] (6.6+\i*1.0,0) circle (0.3);
  \foreach \i in {0,1} \fill[omDef] (9.6+\i*1.0,0) circle (0.3);
  \node[below, font=\small] at (5.6,-0.7)
    {$3 + 2 = 5$：三和二合起来是五};
\end{tikzpicture}

{\small 用图画表示加法：两组合在一起，
合成一组，共五个。}
\end{omfigure}

\begin{example}[加数的位置可以交换]\label{ex:g1:addition:order}
$2 + 3$ 和 $3 + 2$ 都是把同样的弹珠合在一起，结果都是 $5$。
计算 $2 + 9$ 时，从较大的数开始更快：
先说“九”，再往后数两个——“十，十一”。于是
$2 + 9 = 11$。
\end{example}

\section{十的好朋友}

\begin{example}[凑成十]\label{ex:g1:addition:friends}
两个数合起来是 $10$，它们就叫“十的好朋友”：
\[
1 + 9, \quad 2 + 8, \quad 3 + 7, \quad 4 + 6, \quad 5 + 5 .
\]
十根手指可以把它们都表示出来：举起 $3$ 根手指，
那么弯着的 $7$ 根就是它的好朋友。熟记十的好朋友，
很多加法都能算得很快。
\end{example}

\begin{omfigure}
\begin{tikzpicture}[scale=0.72]
  % ten-frame: 2 rows of 5 cells
  \draw[thick] (0,0) grid[step=1] (5,2);
  \draw[very thick] (0,0) rectangle (5,2);
  \foreach \i in {0,1,2} \fill[omDef] (\i+0.5,1.5) circle (0.3);
  \foreach \i in {3,4} \draw[omThm, thick] (\i+0.5,1.5) circle (0.3);
  \foreach \i in {0,...,4} \draw[omThm, thick] (\i+0.5,0.5) circle (0.3);
  \node[right, font=\small, align=left] at (5.6,1.0)
    {一个\emph{十格框}：$10$ 个格子。\\
     已填 $3$ 格，空 $7$ 格：$3 + 7 = 10$};
\end{tikzpicture}

{\small 十格框一眼就能看出每一对十的好朋友：
填色的部分与空的部分合起来总是 $10$。}
\end{omfigure}

\begin{method}[凑十法]\label{met:g1:addition:crossten}
计算 $8 + 5$ 时：
\begin{enumerate}
  \item 从 $5$ 中取出 $8$ 凑成 $10$ 所需的数，即
        $2$（它的好朋友）；
  \item $8 + 2 = 10$，还剩 $3$，它来自 $5$；
  \item $10 + 3 = 13$。所以 $8 + 5 = 13$。
\end{enumerate}
经过 $10$，就把一个难算的和变成两个容易算的和。
\end{method}

\begin{omfigure}
\begin{tikzpicture}[scale=0.62]
  % ten-frame with 8 blue, 2 red (borrowed from the 5)
  \draw[thick] (0,0) grid[step=1] (5,2);
  \draw[very thick] (0,0) rectangle (5,2);
  \foreach \i in {0,...,4} \fill[omDef] (\i+0.5,1.5) circle (0.3);
  \foreach \i in {0,1,2} \fill[omDef] (\i+0.5,0.5) circle (0.3);
  \foreach \i in {3,4} \fill[omThm] (\i+0.5,0.5) circle (0.3);
  % 3 red left outside the frame
  \foreach \i in {0,1,2} \fill[omThm] (6.2+\i,1.0) circle (0.3);
  \node[below, font=\small] at (2.5,-0.25) {十格框满了：$10$};
  \node[below, font=\small] at (7.2,0.55) {多出 $3$ 个};
\end{tikzpicture}

{\small 在十格框上看 $8 + 5$：$8$ 个蓝色计数片借用 $2$ 个
从 $5$ 个红色里取，把框填满——一满十，再多 $3$ 个，
所以是 $13$。}
\end{omfigure}

\begin{omfigure}
\begin{tikzpicture}[scale=0.6]
  \draw[-Stealth] (-0.4,0) -- (15.4,0);
  \foreach \i in {0,...,15} \draw (\i,-0.1) -- (\i,0.1);
  \foreach \i in {0,5,8,10,13,15}
    \node[below, font=\footnotesize] at (\i,-0.12) {\i};
  \draw[-Stealth, omDef, thick] (8,0.3) .. controls (9,0.95) ..
    (10,0.3);
  \node[omDef, font=\small] at (9,1.05) {$+2$};
  \draw[-Stealth, omThm, thick] (10,0.3) .. controls (11.5,1.15) ..
    (13,0.3);
  \node[omThm, font=\small] at (11.5,1.2) {$+3$};
\end{tikzpicture}

{\small 在数轴上计算 $8 + 5$：先跳到 $10$，再跳剩下的。}
\end{omfigure}

\begin{example}[相同数相加]\label{ex:g1:addition:doubles}
这些相同数相加的结果值得熟记：
\[
1+1=2, \quad 2+2=4, \quad 3+3=6, \quad 4+4=8, \quad 5+5=10,
\]
\[
6+6=12, \quad 7+7=14, \quad 8+8=16, \quad 9+9=18, \quad 10+10=20 .
\]
相近的两数相加也很容易：$6 + 7$ 比 $6 + 6$ 多一，所以是 $13$。
\end{example}

\section{小问题}

\begin{example}\label{ex:g1:addition:problem}
“莉娅有 $7$ 张贴纸。她又得到 $6$ 张。现在有多少张？”
得到更多就用加法：$7 + 6 = 13$（凑十法：
$7 + 3 = 10$，再加 $+3$）。莉娅有 $13$ 张贴纸。最后一定要
写一句话作答！
\end{example}

\section{练习}

\begin{exercise}[$\star$]\label{exo:g1:addition:1}
画出弹珠并计算：$4 + 3$；\quad $5 + 2$；\quad
$6 + 4$。
\end{exercise}

\begin{exercise}[$\star$]\label{exo:g1:addition:2}
从较大的数开始算：$3 + 8$；\quad
$2 + 12$；\quad $4 + 9$。
\end{exercise}

\begin{exercise}[$\star$]\label{exo:g1:addition:3}
说出与下列各数凑成十的数：$7$；\ $4$；\ $9$；\ $5$；\ $1$。
\end{exercise}

\begin{exercise}[$\star$]\label{exo:g1:addition:4}
用凑十法计算（见\cref{met:g1:addition:crossten}）：
$9 + 4$；\quad $8 + 7$；\quad $6 + 5$。
\end{exercise}

\begin{exercise}[$\star$]\label{exo:g1:addition:5}
先背一背相同数相加，再用来计算：$7 + 8$；\quad
$5 + 6$；\quad $9 + 8$。
\end{exercise}

\begin{exercise}[$\star$]\label{exo:g1:addition:6}
填空：
\[
6 + \;?\; = 10, \qquad
\;?\; + 5 = 12, \qquad
10 + \;?\; = 17 .
\]
\end{exercise}

\begin{exercise}[$\star$]\label{exo:g1:addition:7}
“车上有 $8$ 个小朋友；$5$ 位大人上车。车上
一共有多少人？”写出加法算式和
答句。
\end{exercise}

\begin{exercise}[$\star$]\label{exo:g1:addition:8}
汤姆有 $6$ 辆红色小车和 $7$ 辆蓝色小车。米娅有 $10$ 辆车。谁
拥有的车更多？
\end{exercise}

\begin{exercise}[$\star\star$]\label{exo:g1:addition:9}
三个骰子分别显示 $4$、$6$ 和 $5$。它们的和是多少？（先把两个骰子
的点数相加——选出能凑成十的那两个！）
\end{exercise}
